\clearpage\thispagestyle{empty}
\pdfbookmark[0]{Quellen zum Ceres-Bericht vom Mai 1802}{ceres180205-sources}
\begin{center}\textbf{Quellen zum Ceres-Bericht vom Mai 1802}\end{center}
\noindent Diese KI-Übertragung in heutiges Deutsch folgt dem deutschen Auszug in Carl Friedrich Gauß’ \href{https://gdz.sub.uni-goettingen.de/id/PPN236007467}{\textit{Werken}, Band~VI} (Göttingen, 1874), S.~211. Franz Xaver von Zach berichtet über Ceres und stellt die Beobachtungen vom Seeberg und aus Greenwich den Rechnungen von Carl Friedrich Gauß gegenüber.
\par\smallskip\noindent Von Zachs \href{https://api.digitale-sammlungen.de/iiif/image/v2/bsb10538597_00474/full/1450,/0/default.jpg}{ursprünglicher Aufsatz}, Nr.~XLIX, erschien im Maiheft 1802 der Zeitschrift \textit{Monatliche Correspondenz zur Beförderung der Erd- und Himmels-Kunde} (Gotha, Becker), Band~V, S.~462--481. Der Auszug in den Werken umfasst Teile der \href{https://api.digitale-sammlungen.de/iiif/image/v2/bsb10538597_00485/full/1450,/0/default.jpg}{S.~473} und \href{https://api.digitale-sammlungen.de/iiif/image/v2/bsb10538597_00486/full/1450,/0/default.jpg}{474}. Die frühere Bemerkung zur Zeitminute der Beobachtung vom 19. Februar steht auf \href{https://api.digitale-sammlungen.de/iiif/image/v2/bsb10538597_00479/full/1450,/0/default.jpg}{S.~467}. Der Aufsatz endet auf \href{https://api.digitale-sammlungen.de/iiif/image/v2/bsb10538597_00493/full/1450,/0/default.jpg}{S.~481}. Die Digitalisate der Zeitschrift stammen von der Bayerischen Staatsbibliothek.
\par\smallskip\noindent Übertragung und Satz wurden mit KI erstellt.
\clearpage
